\section{Artifact Availability and Reproduction}
\label{sec:artifact}

The complete source code, datasets, benchmark harness, and reproducibility
manifests are released under the MIT License at
{\small \url{https://github.com/BlackSamuron0305/semantic-fact-db}}.
Table~\ref{tab:reproducibility} records the exact commit, interpreter,
lock-file hash, dataset checksum, and seed behind every number in this
paper; \texttt{results/paper\_suite\_reproducibility.json} carries the
same manifest alongside the raw benchmark output.

\begin{table}[ht]
\centering
\caption{Reproducibility metadata for the reported run.}
\label{tab:reproducibility}
\begin{tabular}{@{}ll@{}}
\toprule
Property & Value \\
\midrule
Git commit & \texttt{\githash} \\
Python version & \texttt{\pythonversion} \\
uv lock hash & \texttt{\uvlockhash} \\
Dataset checksum & \texttt{\datachecksum} \\
Random seed & \texttt{\sfdbseed} \\
Generation date & \texttt{\generationdate} \\
\bottomrule
\end{tabular}
\end{table}

\label{sec:artifact:reproducibility}
From a clean clone, the following commands reproduce every number and
figure in this paper except the Jena TDB2, Virtuoso, and GraphDB comparisons
and the three LUBM-vs-external-store runs, all four of which require a
separately running external service (see below)
(\texttt{sfdb benchmark} with no \texttt{--size}
argument runs all four scales --- 10 timed runs and 2 warm-up iterations
per query class per Section~\ref{sec:eval:setup} --- and exits non-zero
if any cross-engine verification fails):

\begin{verbatim}
uv sync --group dev
uv run pytest                              # full test suite should pass
uv run sfdb benchmark                      # populates results/ (3 engines)
uv run python scripts/rdflib_reference.py  # external rdflib reference point
uv run python scripts/wikidata_case_study_benchmark.py  # case study
uv run python scripts/lubm_benchmark.py    # LUBM vs. the 3 in-house engines
uv run python scripts/generate_tables.py   # tables + reproducibility macros
uv run python scripts/generate_figures.py  # figures
cd paper && latexmk -pdf main.tex          # main.pdf
\end{verbatim}

The case-study benchmark reads a cached Wikidata snapshot shipped in the
artifact rather than querying Wikidata live. A separate fetch script
(under \texttt{scripts/}) refreshes that snapshot but is an explicit,
optional step, not part of the reproduction sequence above.

The Apache Jena TDB2 comparison (Section~\ref{sec:eval:jena}) requires a
separately running Fuseki server, since it is a genuine external process
rather than code this artifact ships and controls. Requires a Java~17+
runtime, which the reproduction script does not install:

\begin{verbatim}
# One-time: download and start a Fuseki server with a fresh TDB2 dataset
curl -LO https://dlcdn.apache.org/jena/binaries/apache-jena-fuseki-6.2.0.zip
unzip apache-jena-fuseki-6.2.0.zip && cd apache-jena-fuseki-6.2.0
java -jar fuseki-server.jar --loc=./tdb2data --update --port 3030 /sfdb &

# Then, from the repository root:
uv run python scripts/jena_reference.py    # populates results/
uv run python scripts/lubm_external_reference.py --store jena
uv run python scripts/generate_tables.py   # regenerates jena tables too
\end{verbatim}

The OpenLink Virtuoso comparison (Section~\ref{sec:eval:virtuoso})
likewise requires a separately running Virtuoso instance. The simplest
route is the official Docker image:

\begin{verbatim}
# One-time: start a Virtuoso open-source-edition container
docker run -d --name virtuoso-sfdb -p 8890:8890 -p 1111:1111 \
    -e DBA_PASSWORD=sfdb_bench_2026 \
    openlink/virtuoso-opensource-7:latest

# Then, from the repository root:
uv run python scripts/virtuoso_reference.py  # populates results/
uv run python scripts/lubm_external_reference.py --store virtuoso
uv run python scripts/generate_tables.py     # regenerates virtuoso tables too
\end{verbatim}

The Ontotext GraphDB comparison (Section~\ref{sec:eval:graphdb})
likewise requires a separately running GraphDB instance. The simplest
route is the official Docker image; unlike Virtuoso, no authentication
or manual repository setup is needed:

\begin{verbatim}
# One-time: start a GraphDB free-edition container
docker run -d --name graphdb-sfdb -p 7200:7200 ontotext/graphdb:10.7.3

# Then, from the repository root (auto-creates its repository):
uv run python scripts/graphdb_reference.py   # populates results/
uv run python scripts/lubm_external_reference.py --store graphdb
uv run python scripts/generate_tables.py     # regenerates graphdb tables too
\end{verbatim}

To reproduce independent of the host's Python version or
operating system, a \texttt{Dockerfile} at the repository root builds
and runs the test suite in a container:

\begin{verbatim}
docker build -t sfdb .
docker run --rm sfdb                       # runs the test suite
docker run --rm sfdb uv run sfdb benchmark # runs the benchmark suite
\end{verbatim}

The $14$ tests that require a live Jena, Virtuoso, or GraphDB instance
skip inside
the isolated container rather than fail. The image does not bundle TeX
Live or the external stores themselves.

Source code, data generators, and benchmark harnesses are released under
the MIT License.
